\paragraph{Independent verification of the physical kernel in A.4.8.}
The following calculation checks the exact scalar and branches without
using uniqueness for a differential equation.  The Fourier convention is
$\widehat h(\kappa)=\int_{\mathbb R}h(u)e^{i\kappa u}\,du$.
Let $a=\cosh R$, $B=\sinh R$, $R>0$, and initially assume $B>1$.
In the application $a=17$, $B=12\sqrt2$, and $B^{-2}=1/288$.
Write
\[
 K_{p,\kappa}=
 \frac{2^{p-1}\Gamma((p+i\kappa)/2)\Gamma((p-i\kappa)/2)}
      {\Gamma(p)},\qquad
 \tau(u)=\operatorname{arsinh}(B\cosh u).
\]
For $\Re s>0$ set
\[
 f_s(r)=2^{-s}(\cosh r)^{-s}
 {}_2F_1\!\left(\frac s2,\frac s2;s+\frac12;\operatorname{sech}^2r\right).
\]
All powers of positive real numbers use the real logarithm.
The Pfaff transformation, DLMF (15.8.1), with $x=\sinh r>0$,
gives
\[
 f_s(r)=2^{-s}x^{-s}
 {}_2F_1\!\left(\frac s2,\frac{s+1}2;s+\frac12;-x^{-2}\right).
\]
Termwise differentiation therefore yields
\[
 \frac{f_s'(r)}{\cosh r}
 =-s2^{-s}x^{-s-1}
 {}_2F_1\!\left(\frac{s+1}2,\frac{s+2}2;s+\frac12;-x^{-2}\right).
\]
At $r=\tau(u)$ the argument has modulus at most $B^{-2}<1$.
The series is absolutely and uniformly convergent after multiplication
by the integrable majorant $\cosh(u)^{-\Re s-1}$, locally uniformly
in $s$ in the right half-plane.  It may consequently be integrated
term by term against $e^{i\kappa u}$ for real $\kappa$.
The identity
\[
 \frac{K_{s+1+2n,\kappa}}{K_{s+1,\kappa}}
 =\frac{((s+1+i\kappa)/2)_n((s+1-i\kappa)/2)_n}
 {((s+1)/2)_n((s+2)/2)_n}
\]
cancels the two numerator factors of the differentiated series. Hence
\begin{align*}
 \int_{\mathbb R}
 \frac{a f_s'(\tau(u))}{\cosh\tau(u)}e^{i\kappa u}\,du
 &=-sa2^{-s}B^{-s-1}K_{s+1,\kappa}\,\mathcal F_s(\kappa),\\
 \mathcal F_s(\kappa)
 &={}_2F_1\!\left(\frac{s+1+i\kappa}2,
                  \frac{s+1-i\kappa}2;
                  s+\frac12;-B^{-2}\right).
\end{align*}

Independently, the Legendre formula already proved in A.4.7 is
\[
 m_s(\kappa)=\frac{2^{s-1/2}}{\sqrt{\pi B}}
 \Gamma\!\left(\frac{s+1+i\kappa}2\right)
 \Gamma\!\left(\frac{s+1-i\kappa}2\right)
 P_{-1/2+i\kappa}^{\,1/2-s}(a/B).
\]
DLMF (14.3.18), followed by the Pfaff transformation (15.8.1), gives,
for $x>1$,
\[
 P_\nu^{-\mu}(x)=
 \frac{(x^2-1)^{\mu/2}}{2^\mu\Gamma(\mu+1)}
 {}_2F_1\!\left(\frac{\mu-\nu}2,
                  \frac{\mu+\nu+1}2;\mu+1;1-x^2\right).
\]
Here DLMF's normalized hypergeometric function has been converted
using $\mathbf F(A,B;C;z)={}_2F_1(A,B;C;z)/\Gamma(C)$.
Substitute $x=a/B$, $\mu=s-1/2$, and $\nu=-1/2+i\kappa$.
Euler's transformation (also DLMF (15.8.1)) replaces the resulting
first two parameters $(s\pm i\kappa)/2$ by
$(s+1\pm i\kappa)/2$ and introduces precisely
$(1+B^{-2})^{1/2}=a/B$.  Thus
\[
 m_s(\kappa)=
 \frac{a B^{-s-1}}
      {\sqrt\pi\,\Gamma(s+1/2)}
 \Gamma\!\left(\frac{s+1+i\kappa}2\right)
 \Gamma\!\left(\frac{s+1-i\kappa}2\right)
 \mathcal F_s(\kappa).
\]
Comparison with the preceding Fourier integral proves exactly
\[
 \boxed{\quad
 m_s(\kappa)=\int_{\mathbb R}k_s(u)e^{i\kappa u}\,du,
 \qquad
 k_s(u)=-\frac{\Gamma(s)}{\sqrt\pi\,\Gamma(s+1/2)}
       \frac{a f_s'(\tau(u))}{\cosh\tau(u)}.
 \quad}
\]
The original argument $\operatorname{sech}^2\tau$ lies in $(0,1/289]$,
and the transformed one lies in $[-1/288,0)$.
The transformations therefore cross no cut; $a/B$ and
$1+B^{-2}$ are positive.  This checks all branch choices in the
application.  The identity for arbitrary $B>0$, if desired, follows
by analytic continuation in $R>0$ or the differential-equation proof
in A.4.8; that extension is not needed here.

\paragraph{An elementary scalar check.}
At $s=1$,
$f_1(r)=\arctan(e^{-r})$ and $f_1'(r)=-(2\cosh r)^{-1}$.
Since $\Gamma(1)/(\sqrt\pi\Gamma(3/2))=2/\pi$ and
$K_{1,0}=\pi$, the inverse Fourier kernel of $K_{1,0}m_1$ is
\[
 \pi k_1(u)=\frac{a}{1+B^2\cosh^2u}.
\]
This is positive and agrees with direct Fourier inversion at $s=1$.

\paragraph{Uniform outgoing estimate.}
Writing $v_t(r)=\sqrt{\cosh r}\,f_{1/2+it}(r)$ gives
\[
 v_t''+(t^2-V(r))v_t=0,\qquad
 V(r)=\tfrac14\operatorname{sech}^2r,\qquad
 v_t(r)\sim2^{-1/2}e^{-itr}.
\]
The outgoing Volterra equation has kernel bounded in modulus by
$1/t$.  Put $J(r)=\int_r^\infty V(v)\,dv=(1-\tanh r)/4$.
Its successive approximations imply, for $t\geq1$,
\[
 \left|v_t(r)-2^{-1/2}e^{-itr}\right|
 \leq2^{-1/2}(e^{J(r)/t}-1),\qquad
 \left|v_t'(r)+it2^{-1/2}e^{-itr}\right|
 \leq2^{-1/2}J(r)e^{J(r)/t}.
\]
Both are uniform for all $r\geq R_*$, and
$J(r)\leq e^{-2r}/2$.  These inequalities independently verify
the uniform first-order estimates in A.4.8.  They do not estimate
the sum of endpoint Busemann contributions in the relative wave
composition, and no nonvanishing conclusion is drawn from them alone.

\paragraph{Sources checked.}
The full DLMF equations (14.3.18) and (15.8.1), their notation and
principal-branch conventions were opened at
\texttt{https://dlmf.nist.gov/14.3} and
\texttt{https://dlmf.nist.gov/15.8} on 15 November 2026.
